\section{Evaluation Metrics, Benchmarking Testbeds, and Empirical Analysis}
\label{sec:benchmarks}

Rigorous quantification of security postures in agentic systems requires standardized metrics capable of capturing non-deterministic decision trajectories, effector side-effects, and multi-step kill chains~\cite{picek2025security}.

\subsection{Formal Security Metrics Formulation}
\label{subsec:metrics_formulation}

Let $\mathcal{S}$ denote a corpus of benchmark test scenarios, and let $\mathcal{A}$ represent the sequence of discrete actions executed by the agent. We formalize eight foundational security metrics:

\subsubsection{Unsafe Action Rate (UAR)}
The proportion of test scenarios that induce at least one policy-violating action is defined as
\begin{equation}
    \text{UAR} = \frac{1}{|\mathcal{S}|} \sum_{s \in \mathcal{S}} \mathbf{1}\{\text{violation}(s)\},
\end{equation}
where $\mathbf{1}\{\cdot\}$ is the indicator function. Variants of UAR can isolate specific fault modes, such as Tool-Misuse UAR or Exfiltration UAR. Lower values indicate superior security.

\subsubsection{Policy Adherence Rate (PAR)}
The proportion of executed actions that strictly satisfy developer-defined security policies is given by
\begin{equation}
    \text{PAR} = \frac{1}{|\mathcal{A}|} \sum_{a \in \mathcal{A}} \mathbf{1}\{\text{compliant}(a)\}.
\end{equation}
While in single-action adversarial tests PAR approximates $1 - \text{UAR}$, in complex multi-step trajectories PAR and UAR diverge, as an agent may execute multiple benign actions before a single fatal violation.

\subsubsection{Privilege Escalation Distance (PED)}
Defined over the causal threat graph $\mathcal{G}_{\text{threat}} = (\mathcal{V}, \mathcal{E})$, the privilege escalation distance satisfies
\begin{equation}
    \text{PED} = \min_{u \in \mathcal{V}_{\text{source}}, \, \pi \in \mathcal{V}_{\text{asset}}} \text{dist}_{\mathcal{G}}(u, \pi),
\end{equation}
where $\text{dist}_{\mathcal{G}}(u, \pi)$ is the shortest directed path length from an untrusted input node $u$ to a protected asset or privileged effector $\pi$. Higher PED signifies greater defense-in-depth barriers.

\subsubsection{Time-to-Contain (TTC)}
The duration from the initiation of an exploit $t_{\text{start}}$ to successful containment $t_{\text{contain}}$ is given by
\begin{equation}
    \text{TTC} = t_{\text{contain}} - t_{\text{start}}.
\end{equation}
TTC quantifies the latency of runtime anomaly detection and automated kill-switches.

\subsubsection{Patch Half-Life (PHL)}
Let $N(t)$ denote the number of production agent instances remaining vulnerable to a disclosed vulnerability at time $t$. The Patch Half-Life is defined as the minimum duration $\Delta t$ satisfying
\begin{equation}
    N(t_0 + \Delta t) = \frac{1}{2} N(t_0).
\end{equation}
Shorter PHL reflects superior organizational responsiveness across prompt, schema, and dependency patch cycles.

\subsubsection{Retrieval Risk Score (RRS)}
For a retrieved document bundle $\mathcal{B} = \{d_1, d_2, \dots, d_k\}$, let $r(d_i) \in [0, 1]$ denote the assessed risk of document $d_i$ based on provenance, presence of imperative keywords, and anomalous similarity. The aggregate retrieval risk is defined as
\begin{equation}
    \text{RRS}(\mathcal{B}) = \sum_{i=1}^k \alpha_i r(d_i), \quad \text{subject to} \quad \sum_{i=1}^k \alpha_i = 1,
\end{equation}
where $\alpha_i$ represents the model's reliance weight (e.g., attention allocation or citation rank). High RRS triggers human-in-the-loop escalation or effector disabling.

\subsubsection{Out-of-Role Action Rate (OORAR)}
Given an operational role contract $\mathcal{C}$ defining permissible tool verbs and arguments, the out-of-role action rate is
\begin{equation}
    \text{OORAR} = \frac{|\{a \in \mathcal{A} : \neg \text{allowed}_{\mathcal{C}}(a)\}|}{|\mathcal{A}|}.
\end{equation}
Spikes in OORAR signal model drift, prompt confusion, or active adversarial hijacking.

\subsubsection{Cost-Exploit Susceptibility (CES)}
The expected financial loss incurred under an adversarial policy $\pi_{\text{adv}}$ over a finite execution horizon $T$ is given by
\begin{equation}
    \text{CES}_T(\pi_{\text{adv}}) = \mathbb{E}_{\pi_{\text{adv}}} \left[ \sum_{t=1}^T c(a_t) \right],
\end{equation}
where $c(a_t)$ denotes the marginal financial cost (compute fees, third-party API billing, or serverless invocation costs) of action $a_t$.

\subsection{Comparative Analysis of Empirical Benchmarks}
\label{subsec:benchmark_landscape}

Table~\ref{tab:benchmark_comparison} provides a comparative taxonomy of prominent empirical benchmarks developed to evaluate LLM agent security. 

\begin{table}[t]
\centering
\caption{Systematic Comparison of Contemporary Empirical Security Benchmarks and Testbeds for Agentic AI.}
\label{tab:benchmark_comparison}
\footnotesize
\begin{tabularx}{\columnwidth}{l p{2.1cm} p{2.4cm} X X c}
\toprule
\textbf{Benchmark} & \textbf{Target System} & \textbf{Evaluation Environment} & \textbf{Threat Vectors Evaluated} & \textbf{Primary Metrics} & \textbf{Year / Venue} \\
\midrule
\textbf{AgentDojo}~\cite{debenedetti2024agentdojo} & Autonomous Tool Agents & Dynamic simulated web, email, calendar suite & Indirect prompt injection, goal hijacking & Attack Success Rate (ASR), Task Utility & 2024 / NeurIPS \\
\textbf{LLMSmith}~\cite{liu2024demystifying} & LLM Frameworks (LangChain, AutoGPT) & Native OS execution & Remote Code Execution (RCE), command injection & Vulnerability count, Exploit rate & 2024 / ACM CCS \\
\textbf{BIPIA}~\cite{yi2025benchmarking} & QA \& Retrieval LLMs & Static web, email, and document corpora & Indirect prompt injection across diverse formats & ASR, Refusal Rate & 2025 / ACM KDD \\
\textbf{Agent Smith}~\cite{gu2024agent} & Multimodal LLM Swarms & Multi-agent chat network & Self-replicating prompt worms, swarm infection & Infection rate, Propagation velocity & 2024 / ICML \\
\textbf{PoisonedRAG}~\cite{zou2025poisonedrag} & RAG Systems & Vector DBs (Dense \& Sparse retrievers) & Knowledge base poisoning, targeted corruption & Targeted Misinformation Rate & 2025 / USENIX Sec \\
\textbf{Blocker RAG}~\cite{shafran2025machine} & Dense Retrieval RAG & Enterprise documentation corpora & Denial of retrieval, blocker document jamming & Top-$k$ Jamming Rate, Utility Drop & 2025 / USENIX Sec \\
\textbf{AgentPoison}~\cite{chen2024agentpoison} & Long-term Memory Agents & Vector databases \& episodic memory stores & Memory poisoning, backdoor triggers & Backdoor ASR, Task Accuracy & 2024 / NeurIPS \\
\textbf{InjecAgent}~\cite{zhan2024injecagent} & ReAct Tool Agents & Synthetic tool environments & Tool parameter hijacking, data exfiltration & ReAct Step Hijack Rate & 2024 / ACL Findings \\
\textbf{ToolEmu}~\cite{ruan2023identifying} & Function Calling Agents & LM-emulated tool execution sandbox & Unintended real-world tool side-effects & Severe Hazard Rate & 2024 / ICLR \\
\bottomrule
\end{tabularx}
\end{table}


Early benchmarks such as BIPIA~\cite{yi2025benchmarking} pioneered the systematic measurement of indirect prompt injection across diverse text formats; however, their evaluation was confined to static question-answering environments. In contrast, modern testbeds emphasize interactive tool execution. \textbf{AgentDojo}~\cite{debenedetti2024agentdojo} provides a realistic, dynamic environment featuring simulated email, web, and calendar services, exposing how indirect prompt injection causes agents to compromise user data. 

In the software execution domain, \textbf{LLMSmith}~\cite{liu2024demystifying} evaluated 11 major open-source agent frameworks, uncovering 19 critical RCE vulnerabilities arising from unchecked tool parameter handling. For multi-agent systems, \textbf{Agent Smith}~\cite{gu2024agent} established the first rigorous benchmark for measuring the velocity and infection rate of self-replicating prompt worms across interconnected multimodal agents.

\subsection{Evaluation Pitfalls and Testbed Rigor}
\label{subsec:eval_pitfalls}

Conducting reproducible security evaluations on agentic AI presents significant methodological pitfalls:
\begin{itemize}
    \item \textbf{Evaluator Leakage and LLM-as-a-Judge Bias:} Employing the model under test (or a model with shared pre-training weights) to evaluate trajectory safety introduces severe bias. Evaluation harnesses must employ deterministic, side-channel criteria, such as AST parsing of tool parameters, filesystem diff tracking, and network egress packet captures.
    \item \textbf{Non-Determinism and State Carry-Over:} Because LLM decoding is stochastic, single-run evaluations yield high variance. Reliable evaluation requires multiple randomized trials with time-stratified sampling and strict between-trial state resets (flushing vector caches, resetting virtual filesystems, and clearing conversation memories)~\cite{dehghantanha2026sok}.
    \item \textbf{Continuous Red-Teaming Integration:} Because foundation models, system prompts, and tool libraries undergo frequent updates, security benchmarks must be embedded directly within CI/CD pipelines with automated regression gates (e.g., failing builds if UAR increases by more than 2 percentage points).
\end{itemize}
